\documentclass[12pt]{article}
\usepackage[utf8]{inputenc}
\usepackage[T1]{fontenc}
\usepackage{hyperref}
\title{Marriammu Language Resources for K-12 Curriculum Development}
\author{Compiled by Claude Sonnet 5.5 (claude-sonnet-5-5) for Andr\'as Kornai}
\date{}
\begin{document}
\maketitle

\section{Introduction}
Marriammu (ISO 639-3 xru; also Marri Amu, Mareammu) is a critically endangered Western Daly language of the Northern Territory, Australia. Sources found in search treat it as a dialect of, or closely tied to, Marrithiyel (ISO mfr), and Marrithiyel had about 15 speakers in the 2016 census. In some communities no one has grown up speaking the language, but people still identify with it and are working to revitalize it. This search found no school curriculum or teaching materials specific to Marriammu.

\section{Primary Educational Resources}
No primers, textbooks or school programmes for Marriammu turned up in search. Marri Amu was named one of 20 languages prioritised under the Priority Languages Support Project run by First Languages Australia. Any curriculum work would have to start from community consultation and the archival material below.

\section{Linguistic and Reference Materials}
Glottolog lists sources on Marrithiyel and its dialects. These include Tryon's Daly River surveys, Stanner's 1938--39 notes and Ian Green's 1989 ANU PhD thesis on Marrithiyel. Glottolog entry: \url{https://glottolog.org/resource/languoid/id/mari1420}. Overview of the related Marrithiyel language: \url{https://en.wikipedia.org/wiki/Marrithiyel_language}.

\section{Religious and Community Resources}
Joshua Project has a listing for xru: \url{https://joshuaproject.net/languages/xru}. The Mariamo (Marri Ammu) people are described at \url{https://www.Wikipedia.org/wiki/Mariamo}. No scripture translations were found.

\section{Digital and Audio Resources}
The Endangered Languages Project has an entry: \url{https://endangeredlanguages.com/lang/2329}. No audio or video resources were found in search.

\section{Curriculum Development Context}
The language is very severely endangered and the documentation is thin. Any K-12 work would be a heritage and revitalization effort led by the community, probably together with language centres and First Languages Australia. Archival sources such as Stanner's notes and Green's thesis would need to be located and used with the community's permission.
\end{document}